\subsection{Common-reference cancellation at the original optimum}\label{sec:centered60}
The width of an absolute value box can be dominated by state variation common to all actions. Subtracting the same reference function from both Bellman endpoints removes that variation without changing the decision problem. This is a use of the original policy-relevant error principle, not a replacement of the Bellman objective by a zero-policy comparison.

Fix any feasible complete reference $\pi^0$ and write $h_t=J_t^{\pi^0}$. Define its true reference advantage by
\[
 A_t^0(x,a)=c_t(x,a)+\beta_tP_t^a h_{t+1}(x)-h_t(x).
\]
Set $e_t^*(x)=J_t^*(x)-h_t(x)$ and $e_t^\pi(x)=J_t^\pi(x)-h_t(x)$. Their exact recursions are
\begin{align*}
 e_t^*(x)&=\inf_{a\in A_t(x)}
       \{A_t^0(x,a)+\beta_tP_t^a e_{t+1}^*(x)\},\\
 e_t^\pi(x)&=A_t^0(x,\pi_t(x))
             +\beta_tP_t^{\pi_t(x)}e_{t+1}^\pi(x),
 \qquad e_T^*=e_T^\pi=0.
\end{align*}
The reference need not be fitted, and a numerical routine need not evaluate $h_t$ separately if it encloses its signed advantages directly from the primitives.

\begin{corollary}[Reference-centered Bellman certificate]\label{cor:centered60}
Apply the lower-cover and robust upper-candidate construction of Theorem~\ref{thm:bellman-bracket60} to the excess recursions above, using verified enclosures of $A_t^0$ and the original probability kernels. Denote the resulting piecewise excess bounds by $\ell_t$ and $u_t$. Then, at every true state,
\[
 \ell_t(x)\le J_t^*(x)-h_t(x)
       \le J_t^\pi(x)-h_t(x)\le u_t(x).
\]
Consequently $\sup_x\{u_t(x)-\ell_t(x)\}$ is an all-state loss certificate for the original constructed policy against the original continuous-action Bellman optimum. There is no separate reference-approximation allowance when the advantage enclosure is obtained directly from the exact reference law.

For an independently fixed feasible policy $\widehat\pi$, an upper-only excess recursion using its full acquired-actor action range gives $\widehat u_t\ge J_t^{\widehat\pi}-h_t$. The quantity $\sup_x\{\widehat u_t(x)-\ell_t(x)\}$ therefore certifies that particular returned policy. It does not require the policy to be regenerated on the verification grid.
\end{corollary}
\begin{proof}
Subtract $h_t$ from the Bellman equation and from the policy-value equation to obtain the two recursions. The induction in Theorem~\ref{thm:bellman-bracket60} applies verbatim to these transformed equations: kernel monotonicity and the complete feasible-action cover give the lower excess bound, and implementability and the upper continuation bound give the upper excess bound. Both endpoints at a given true state contain the identical $h_t(x)$, which cancels in their difference. The same upper induction applies to any fixed feasible actor when its complete action range is enclosed on every verification cell. Subtracting the lower optimum bound proves the final assertion.
\end{proof}

In the two-date investment economy, choose the original zero-investment policy as reference. Its final-date advantage is the original analytic final-cost difference. At the initial date, use a common innovation to propagate the state difference caused by the proposed action and integrate the difference of subsequent zero-policy costs. Centered polynomial identities evaluate that difference; the final innovation's common variance cancels, while its correct marginal law is retained. The construction therefore does not assume a policy-value oracle, a differentiable acquired policy, or a learned invariant.

The centered experiment is explicitly subsequent to the uncentered one. It retains the unsuccessful original tolerance tests and applies a separately frozen five-rung verification design. The same centered lower table can also certify the previously returned learned and conventional policies. Such revalidation has its own cost and does not retroactively alter the original prospective stopping criterion or its measured clock. This connects the original Bellman-accuracy question to the actual returned objects, rather than only to an alternative constructed policy.
